% ============================================================
%  Appendix A -- Implementation Details and Configuration
% ============================================================
\chapter{Implementation Details}
\label{app:implementation}

This appendix provides implementation details, hyperparameter configurations,
and dataset statistics that complement the main text.

% ============================================================
\section{Codebase Structure}
\label{app:codebase}

The complete implementation is organised into four top-level packages:

\paragraph{\texttt{pcml/}.}
Core library containing:
\begin{itemize}
  \item \texttt{pcml/adapters/} --- Framework adapters (PCGCv2, G-PCC,
    pcc\_geo\_cnn\_v2). Each adapter wraps the corresponding codec via subprocess
    isolation, so each codec runs in its own conda environment with compatible
    dependencies. The \texttt{BaseAdapter} interface exposes
    \texttt{encode()}, \texttt{decode()}, and \texttt{get\_bpp()} methods.

  \item \texttt{pcml/metrics/} --- Quality metrics (\texttt{quality.py}: D1-PSNR,
    D2-PSNR, Chamfer Distance using \texttt{scipy.spatial.cKDTree};
    \texttt{curvature.py}: PCA-based curvature estimation, curvature histograms,
    stratified PSNR).

  \item \texttt{pcml/data/} --- Dataset loaders for PLY files and decoded cloud
    manifests.

  \item \texttt{pcml/visualization/} --- Rate-distortion curve plotting, spatial
    error visualisation.
\end{itemize}

\paragraph{\texttt{postprocessing/}.}
Post-processing model:
\begin{itemize}
  \item \texttt{model.py} --- \texttt{FiLMHeadSparseUNetV2} architecture
    (final model), plus earlier \texttt{FiLMSparseUNet}, \texttt{GatedSparseUNet},
    \texttt{ThresholdGatedUNet} (retained for reference).
  \item \texttt{losses.py} --- \texttt{dynamic\_chamfer\_loss} (production),
    plus earlier \texttt{L1DisplacementLoss}, \texttt{CurvatureWeightedL1},
    \texttt{MinKLoss}, \texttt{StratifiedLoss} (historical).
  \item \texttt{dataset.py} --- \texttt{MultiFrameDataset}: reads decoded cloud
    manifests, applies 48-element cube symmetry augmentation.
  \item \texttt{train.py} --- Training loop with AdamW + warm-up + cosine
    annealing. Supports multi-rate joint training with configurable weights.
  \item \texttt{inference.py} --- Sliding-window inference over full clouds.
    Handles patch overlap, aggregates displacements by averaging.
\end{itemize}

\paragraph{\texttt{scripts/}.}
Organised by function:
\begin{itemize}
  \item \texttt{scripts/setup/} --- Environment setup, dataset symlinks,
    verification.
  \item \texttt{scripts/analysis/} --- Curvature histograms, displacement
    distribution analysis, oversmoothing characterisation.
  \item \texttt{scripts/benchmark/} --- Single-frame and multi-frame codec
    benchmarking.
  \item \texttt{scripts/evaluation/} --- Full-dataset evaluation pipeline
    (\texttt{eval\_full\_dataset.py}, \texttt{compute\_all\_metrics.py},
    \texttt{run\_full\_pipeline.sh}).
  \item \texttt{scripts/plotting/} --- Figure generation for the paper and thesis.
\end{itemize}

\paragraph{\texttt{configs/}.}
YAML configuration files for all training runs. Production config:
\texttt{configs/cls\_v1\_r2.yaml} (v10-nocurv, r2+r4+r6 joint training).

% ============================================================
\section{PCGCv2 Codec Configuration}
\label{app:codec-config}

PCGCv2 is run in its default multi-scale configuration.
Table~\ref{tab:codec-config} lists the compression rates used throughout
this work, with the corresponding quantisation step sizes and measured
bitrates on the 8iVFB dataset.

\begin{table}[h]
  \centering
  \caption{%
    PCGCv2 compression rates and measured bitrates (bpp) on the 8iVFB
    dataset (average over all four sequences and 100 frames each).
  }
  \label{tab:codec-config}
  \small
  \begin{tabular}{ccc}
    \toprule
    \textbf{Rate Index} & \textbf{Quantisation Step} & \textbf{Avg BPP} \\
    \midrule
    r1 & 1 & $\approx 0.025$ \\
    r2 & 2 & $\approx 0.050$ \\
    r3 & 3 & $\approx 0.080$ \\
    r4 & 4 & $\approx 0.150$ \\
    r5 & 5 & $\approx 0.200$ \\
    r6 & 6 & $\approx 0.300$ \\
    r7 & 7 & $\approx 0.450$ \\
    \bottomrule
  \end{tabular}
\end{table}

Note that BPP values vary with sequence and frame due to differences in
point cloud size and geometric complexity.
The values above are approximate averages.
The FiLM head uses the actual per-cloud BPP as the conditioning scalar
at both training and inference time, not the nominal rate index.

% ============================================================
\section{Training Hyperparameters}
\label{app:hyperparams}

Table~\ref{tab:hyperparams} lists the complete set of hyperparameters used
to train the final v10-nocurv model.

\begin{table}[h]
  \centering
  \caption{Complete training hyperparameter configuration (v10-nocurv model).}
  \label{tab:hyperparams}
  \small
  \begin{tabular}{lll}
    \toprule
    \textbf{Category} & \textbf{Hyperparameter} & \textbf{Value} \\
    \midrule
    \multirow{4}{*}{Data} &
      Training sequences & longdress, loot, soldier \\
    & Validation sequence & redandblack \\
    & Frame stride & 12 (every 12th frame) \\
    & Training clouds per rate & $\approx 75$ (3 seq. $\times$ 25 frames) \\
    \midrule
    \multirow{4}{*}{Patches} &
      Patch size & $64^3$ voxels \\
    & Patch stride & 32 (50\% overlap) \\
    & Min points per patch & 10 \\
    & Augmentation & 48 cube symmetries \\
    \midrule
    \multirow{5}{*}{Optimiser} &
      Algorithm & AdamW \\
    & Learning rate ($\eta$) & $10^{-3}$ \\
    & Weight decay ($\lambda$) & $10^{-4}$ \\
    & Momentum ($\beta_1, \beta_2$) & 0.9, 0.999 \\
    & Gradient clipping & None \\
    \midrule
    \multirow{4}{*}{Schedule} &
      Total epochs & 30 \\
    & Warm-up epochs & 3 (linear, $10^{-5} \to 10^{-3}$) \\
    & Annealing epochs & 27 (cosine, $10^{-3} \to 10^{-6}$) \\
    & Best checkpoint & Epoch 24 (min validation loss) \\
    \midrule
    \multirow{4}{*}{Loss} &
      Loss function & Symmetric Chamfer Distance \\
    & Boundary padding & 0 (no padding) \\
    & Rate weights & $[w_{r2}, w_{r4}, w_{r6}] = [1.0, 0.3, 0.05]$ \\
    & Batch size & 8 patches \\
    \midrule
    \multirow{4}{*}{Architecture} &
      Input channels & 3 (normalised $xyz$) \\
    & Backbone channels & 32 / 64 / 128 / 64 / 32 \\
    & FiLM MLP & Linear(1,64) + ReLU + Linear(64,64) \\
    & Output & $5 \cdot \tanh(\Delta) \in [-5,5]^3$ \\
    \midrule
    \multirow{3}{*}{Hardware} &
      GPU & NVIDIA RTX 3090 (24 GB VRAM) \\
    & CPU & 12 cores / 24 threads \\
    & RAM & 64 GB \\
    \bottomrule
  \end{tabular}
\end{table}

% ============================================================
\section{8iVFB Dataset Statistics}
\label{app:dataset}

The 8iVFB~\cite{8ivfb} (8i Voxelized Full Bodies) dataset contains
four volumetric video sequences of human subjects, captured at 10 fps
and voxelized at resolution $G = 1024$:
\textit{longdress}, \textit{loot}, \textit{redandblack}, and \textit{soldier}.

Table~\ref{tab:dataset-stats} summarises the dataset statistics.

\begin{table}[h]
  \centering
  \caption{%
    8iVFB dataset statistics (averages over the first 100 frames per sequence
    at the original resolution $G = 1024$).
  }
  \label{tab:dataset-stats}
  \small
  \begin{tabular}{lcccc}
    \toprule
    \textbf{Sequence} & \textbf{Frames} &
    \textbf{Avg Points} & \textbf{Avg High-$\kappa$\%} & \textbf{Subject} \\
    \midrule
    longdress   & 300 & $\approx 857\,000$ & $\approx 33\%$ & Female, dress \\
    loot        & 300 & $\approx 805\,000$ & $\approx 33\%$ & Male, casual \\
    redandblack & 300 & $\approx 757\,000$ & $\approx 33\%$ & Male, T-shirt \\
    soldier     & 300 & $\approx 1\,089\,000$ & $\approx 33\%$ & Male, uniform \\
    \bottomrule
  \end{tabular}
\end{table}

The high-curvature fraction (top third of points by PCA curvature $\kappa$)
is approximately 33\% by construction of the stratification (tertile split).
The dataset was licensed from 8i Labs and is available at the MPEG content
repository.

For training, we use frames 1--300 of \textit{longdress}, \textit{loot},
and \textit{soldier}, with frame stride 12 (25 frames per sequence).
\textit{redandblack} is held out entirely for validation and evaluation.
For the full 400-frame evaluation, we use 100 frames from each sequence
(sampled with stride 3 from frames 1--300).

% ============================================================
\section{Evaluation Pipeline}
\label{app:eval-pipeline}

The evaluation pipeline consists of three stages:

\paragraph{Stage 1: Inference.}
\texttt{eval\_full\_dataset.py --infer\_only} runs the post-processing model
on all 400 frames $\times$ 3 rates using GPU-accelerated sliding-window
inference.
Refined clouds are saved as \texttt{.npy} files
(format: $N \times 3$ float32).
With 399 cached frames and 1 new inference, this stage takes $\approx 8$ seconds.
For a full cold run (no cache), the estimated time is $\approx 30$ minutes.

\paragraph{Stage 2: Metric computation.}
\texttt{compute\_all\_metrics.py} computes D1-PSNR, flat D1-PSNR, and edge
D1-PSNR for each frame.
Key optimisations:
\begin{itemize}
  \item \texttt{cKDTree(original)} is built once per frame and reused for all
    three metric types (forward distances), saving $\approx 3\times$ tree
    construction time.
  \item \texttt{rev\_sq\_dists} (backward distances, from original to decoded/refined)
    are computed once for D1 and reused for the stratified metric.
  \item \texttt{workers=-1} enables multi-core parallelism across frames.
\end{itemize}
Average throughput: $\approx 3$ seconds per frame.
Total for 400 frames: $\approx 20$ minutes per rate.

\paragraph{Stage 3: Aggregation.}
\texttt{run\_full\_pipeline.sh} orchestrates Stages 1--2 in parallel across
the 3 rates and merges the per-rate CSVs into a single
\texttt{nocurv\_ep24\_all.csv} with 1200 rows.

% ============================================================
\section{Software Dependencies}
\label{app:dependencies}

Table~\ref{tab:dependencies} lists the key software dependencies.

\begin{table}[h]
  \centering
  \caption{Key software dependencies for the \texttt{pcgcv2} conda environment.}
  \label{tab:dependencies}
  \small
  \begin{tabular}{lll}
    \toprule
    \textbf{Package} & \textbf{Version} & \textbf{Purpose} \\
    \midrule
    Python & 3.8 & Runtime \\
    PyTorch & 1.10.1+cu113 & Deep learning framework \\
    MinkowskiEngine & 0.5.4 & Sparse convolution \\
    CUDA & 11.3 & GPU acceleration \\
    GCC/G++ & 10.x & MinkowskiEngine compilation \\
    NumPy & 1.21 & Array operations \\
    SciPy & 1.7 & cKDTree for nearest-neighbour \\
    Open3D & 0.13 & Point cloud I/O \\
    PyYAML & 5.4 & Configuration files \\
    tqdm & 4.62 & Progress bars \\
    pandas & 1.3 & CSV result aggregation \\
    \midrule
    PCGCv2 & --- & Learned codec (git submodule) \\
    mpeg-pcc-tmc13 & --- & G-PCC codec (git submodule) \\
    pcc\_geo\_cnn\_v2 & --- & Alternative codec (git submodule) \\
    \bottomrule
  \end{tabular}
\end{table}

MinkowskiEngine must be built from source with CUDA 11.3 and GCC $\leq$ 10.
The \texttt{scripts/setup/setup\_conda\_env.sh} script automates this process,
including detecting the CUDA version and selecting the correct GCC.

% ============================================================
\section{Displacement Distribution at Each Rate}
\label{app:displacement}

Table~\ref{tab:displacement-full} provides the full displacement statistics
at all seven PCGCv2 rates for a representative 8iVFB sequence (\textit{redandblack}).
These statistics motivate the multi-rate training strategy and the static
rate weights.

\begin{table}[h]
  \centering
  \caption{%
    Displacement statistics (GT nearest-neighbour displacement magnitude
    $\|d_{gt}(\hat{p}_i)\|$ in voxels) at each PCGCv2 rate for the
    \textit{redandblack} sequence, averaged over 10 representative frames.
    ``Zero fraction'' is the fraction of points with $\|d_{gt}\| = 0$.
  }
  \label{tab:displacement-full}
  \small
  \begin{tabular}{cccc}
    \toprule
    \textbf{Rate} & \textbf{Mean $\|d_{gt}\|$} &
    \textbf{90th pct.} & \textbf{Zero fraction} \\
    \midrule
    r1 & 0.32 & 1.41 & 40\% \\
    r2 & 0.15 & 1.06 & 60\% \\
    r3 & 0.09 & 0.82 & 71\% \\
    r4 & 0.06 & 0.62 & 79\% \\
    r5 & 0.05 & 0.50 & 83\% \\
    r6 & 0.04 & 0.38 & 87\% \\
    r7 & 0.02 & 0.21 & 89\% \\
    \bottomrule
  \end{tabular}
\end{table}

The table shows clearly why training at r7 is unproductive: 89\% zero targets
with mean displacement 0.02 voxels (below quantization resolution) leave
essentially no learnable signal.
The three training rates r2, r4, r6 span the range where correction is both
needed (non-trivial zero fraction) and feasible (sub-voxel mean displacement
within the model's $\pm5$-voxel output bound).
